% This file is a chapter fragment, in the same style as the other files under
% version/detailed/.  It is included by version/detailed.tex and must not
% contain a preamble of its own.
%
% Contents: section 2.6 of the curve estimates -- the multiplication-by-n
% endomorphism [n] of an abelian variety, its n-torsion subgroup scheme A[n],
% and the divisibility of the group of rational points.  This is the LaTeX
% source for the three propositions that until now existed only inside
% blueprint/curve_estimates.pdf, together with a proof-level account of the
% route to etaleness of [n] in the good characteristic.

\chapter{The multiplication-by-$n$ isogeny of an abelian variety}
\label{chap:av-isogeny}

The goal of this chapter is \Cref{prop:av-isogeny}: for an abelian variety $A$
of dimension $g$ over a field $K$ and an integer $n\neq0$, the
multiplication-by-$n$ endomorphism $[n]\colon A\to A$ is an isogeny, that is, it
is finite, flat and surjective, and its degree is $n^{2g}$.  Two consequences
are drawn from it: \Cref{prop:av-etale-torsion}, that for $n$ invertible in $K$
the $n$-torsion subgroup scheme $A[n]=\ker[n]$ is finite \'etale of rank
$n^{2g}$, so that over an algebraically closed field
$A[n](K)\cong(\Z/n)^{2g}$; and \Cref{prop:av-divisibility}, that over an
algebraically closed field the group $A(K)$ of rational points is
$n$-divisible.  These are Propositions 2.7, 2.8 and 2.9 of \S2.6 of the curve
estimates, and the labels \texttt{prop:av-isogeny},
\texttt{prop:av-etale-torsion} and \texttt{prop:av-divisibility} are the ones
already used there.

The chapter is written for a reader who has to formalise from it, and is
organised so that such a reader needs to consult nothing else.
\Cref{sec:av-conventions} fixes the conventions and lists the facts that are
already available in Lean, each with the name of the declaration that provides
it.  \Cref{sec:av-three-propositions} states the three propositions and gives
the proofs of \Cref{prop:av-etale-torsion} and \Cref{prop:av-divisibility} from
\Cref{prop:av-isogeny}; those two proofs are complete in Lean.
\Cref{sec:av-etale-route} is the heart: a complete, step-by-step proof that
$[n]$ is \'etale, and then finite, flat and surjective, whenever $n$ is
invertible in $K$, using only deformation theory and the smoothness of $A$, and
in particular using nothing about line bundles.  This is the part that has
since been formalised; it is written so that every step there is a labelled
statement with its own proof, with no step left as ``standard'', and the status
note at the head of that section records where the Lean proof departs from the
decomposition written here.  \Cref{sec:av-cube} treats what the \'etale route
does not give --- the degree $n^{2g}$, and the case where $n$ is divisible by
the characteristic --- in the form of five inputs from the theory of line
bundles on abelian varieties, stated with proof sketches and citations,
each marked as an open obligation because Mathlib has none of the theory they
need.  \Cref{sec:av-status} records the formalisation status of every statement
of the chapter.

\section{Conventions and standing facts}\label{sec:av-conventions}

Throughout this chapter $K$ is a field.  A \emph{group scheme over $K$} is an
object $G$ of the cartesian monoidal category of $K$-schemes together with a
group-object structure on it; we write $\mu\colon G\times_KG\to G$ for the
multiplication, $e\colon\Spec K\to G$ for the identity section and
$\iota\colon G\to G$ for the inverse.  For a $K$-scheme $T$ we write
$G(T):=\Hom_K(T,G)$ for the group of $T$-points, with multiplication
$u\cdot v:=\mu\circ(u,v)$; for a $K$-algebra $R$ we abbreviate
$G(R):=G(\Spec R)$.  All the arguments below are carried out on points, that
is, in the groups $G(T)$; this is also how they are carried out in Lean.

Composition is written in the classical order: $g\circ f$ means ``first $f$,
then $g$''.  The Lean sources use the categorical order and write the same
morphism as \texttt{f} $\gg$ \texttt{g}, so every composite below appears
reversed there.

\begin{definition}[Abelian variety]\label{def:av-abelian-variety}
\lean{AlgebraicGeometry.IsAbelianVariety}
\leanok
An \emph{abelian variety} over $K$ is a group scheme $A$ over $K$ whose
structure morphism $A\to\Spec K$ is proper and geometrically integral.
\end{definition}

The dimension of an abelian variety is deliberately not part of the definition.
Instead, whenever a dimension $g$ is needed it is carried as the hypothesis
that $A\to\Spec K$ is smooth of relative dimension $g$; this is legitimate by
\Cref{prop:av-basic}, and it is how the Lean statements are phrased.

\begin{proposition}[An abelian variety is integral, commutative and smooth]
\label{prop:av-basic}
\uses{def:av-abelian-variety}
\lean{AlgebraicGeometry.isCommMonObj_of_isProper_of_geometricallyIntegral,
AlgebraicGeometry.smooth_of_grpObj}
\leanok
Let $A$ be an abelian variety over $K$.  Then the scheme $A$ is integral, in
particular irreducible, connected and non-empty; the structure morphism
$A\to\Spec K$ is separated, quasi-compact and locally of finite type; the group
scheme $A$ is commutative; and $A\to\Spec K$ is smooth, hence flat and locally
of finite presentation.
\end{proposition}
\begin{proof}
\leanok
Integrality is geometric integrality applied to the trivial field extension,
and the listed properties of $A\to\Spec K$ follow from properness and from
finite type.  Commutativity is \stacksproject{0BFD}.  Smoothness holds because
a geometrically reduced group scheme locally of finite type over a field is
smooth.  All of this is already in Mathlib and in the project's mirror file
\texttt{WeilConjectures/Mathlib/AlgebraicGeometry/Group/Abelian.lean}, which
registers each item as an instance.
\end{proof}

\begin{definition}[The multiplication-by-$n$ endomorphism]\label{def:av-mulby}
\uses{def:av-abelian-variety}
\lean{AlgebraicGeometry.IsAbelianVariety.mulBy, CategoryTheory.GrpObj.npow}
\leanok
Let $G$ be a group scheme over $K$ and $n\in\N$.  The
\emph{multiplication-by-$n$ endomorphism} $[n]\colon G\to G$ is the $n$-th
power of the identity in the group $G(G)=\Hom_K(G,G)$, that is, the unique
morphism whose effect on points is $x\mapsto x^n$.
\end{definition}

\begin{lemma}[Basic algebra of $[n]$]\label{lem:av-mulby-points}
\uses{def:av-mulby}
\lean{CategoryTheory.GrpObj.comp_npow, CategoryTheory.GrpObj.npow_comp,
CategoryTheory.GrpObj.npow_comp_npow}
\leanok
Let $G$ be a group scheme over $K$, let $T$ be a $K$-scheme and let
$u\in G(T)$.  Then $[n]\circ u=u^n$ in $G(T)$.  Consequently $[0]$ is the
constant morphism at the identity, $[1]=\mathrm{id}$, $[m]\circ[n]=[mn]$, and
if $f\colon G\to H$ is a homomorphism of group schemes then $f\circ[n]=[n]\circ
f$.  If $G$ is commutative then $[n]$ is itself a homomorphism of group
schemes.
\end{lemma}
\begin{proof}
\leanok
The identity $[n]\circ u=u^n$ is the defining property of $[n]$: precomposition
with a fixed $u$ is the group homomorphism $G(G)\to G(T)$ induced by $u$, and
it sends $\mathrm{id}$ to $u$, hence $\mathrm{id}^n$ to $u^n$.  Every remaining
assertion is the corresponding identity in the abstract groups $G(T)$, applied
for all $T$ simultaneously and then transported back by Yoneda; for the last
one, $x\mapsto x^n$ is a group homomorphism exactly when the group is
commutative.
\end{proof}

\begin{definition}[The $n$-torsion subgroup scheme]\label{def:av-torsion}
\uses{def:av-mulby, prop:av-basic}
\lean{AlgebraicGeometry.IsAbelianVariety.torsion}
\leanok
Let $A$ be an abelian variety over $K$ and $n\in\N$.  Since $A$ is commutative,
$[n]\colon A\to A$ is a homomorphism of group schemes, and its \emph{kernel}
\[
  A[n]:=\ker[n]:=A\times_{[n],\,A,\,e}\Spec K
\]
is again a group scheme over $K$, the \emph{$n$-torsion subgroup scheme} of
$A$.
\end{definition}

\begin{lemma}[The defining square of $A[n]$]\label{lem:av-torsion-pullback}
\uses{def:av-torsion}
\lean{AlgebraicGeometry.IsAbelianVariety.isPullback_torsion,
AlgebraicGeometry.IsAbelianVariety.torsionHomMulEquiv}
\leanok
The square
\[
  \begin{CD}
    A[n]      @>>>     A\\
    @VVV               @VV{[n]}V\\
    \Spec K   @>{e}>>  A
  \end{CD}
\]
is cartesian.  In particular $A[n]\to\Spec K$ is the base change of $[n]$ along
the identity section, so every property of morphisms of schemes that is stable
under base change passes from $[n]$ to $A[n]\to\Spec K$.  On points, for every
$K$-scheme $T$ the group $A[n](T)$ is the subgroup $\{u\in A(T):u^n=1\}$ of
$A(T)$; taking $T=\Spec K$, the group $A[n](K)$ is the $n$-torsion of $A(K)$.
\end{lemma}
\begin{proof}
\leanok
The square is cartesian by the definition of the kernel as a fibre product.
The description on points is the universal property of that fibre product
combined with $\Hom_K(T,\Spec K)$ being a singleton: a $T$-point of $A[n]$ is a
$T$-point $u$ of $A$ with $[n]\circ u=1$, and $[n]\circ u=u^n$ by
\Cref{lem:av-mulby-points}.
\end{proof}

\begin{definition}[Degree of a finite flat morphism]\label{def:av-degree}
\lean{AlgebraicGeometry.Scheme.Hom.finrank}
\leanok
Let $f\colon X\to Y$ be a morphism of schemes and $y\in Y$.  When $f$ is finite
and flat, $f_*\O_X$ is locally free of finite rank, and we write $\deg_yf$ for
the rank of $(f_*\O_X)_y$ over $\O_{Y,y}$.  If $Y$ is connected this is
independent of $y$ and is written $\deg f$.
\end{definition}

\begin{lemma}[The degree is invariant under base change]
\label{lem:av-degree-base-change}
\uses{def:av-degree}
\lean{AlgebraicGeometry.Scheme.Hom.finrank_of_isPullback,
AlgebraicGeometry.Scheme.Hom.one_le_finrank_iff_surjective}
\leanok
Let
\[
  \begin{CD}
    X'     @>>>   X\\
    @V{f'}VV      @VV{f}V\\
    Y'     @>{g}>> Y
  \end{CD}
\]
be a cartesian square with $f$ finite and flat.  Then $f'$ is finite and flat
and $\deg_{y'}f'=\deg_{g(y')}f$ for every $y'\in Y'$.  Moreover a finite flat
morphism is surjective if and only if its degree is everywhere at least $1$.
\end{lemma}
\begin{proof}
\leanok
Finiteness and flatness are stable under base change, and the formation of
$f_*\O_X$ commutes with flat base change, so this is the invariance of the rank
of a finitely generated projective module under base change.  Both statements
are in Mathlib.
\end{proof}

\begin{theorem}[Zariski's main theorem, finiteness form]
\label{thm:av-zariski-main}
\lean{AlgebraicGeometry.IsFinite.of_isProper_of_locallyQuasiFinite}
\leanok
A proper, locally quasi-finite morphism of schemes is finite.
\end{theorem}
\begin{proof}
\leanok
This is \stacksproject{02LS}; it is in Mathlib.
\end{proof}

\begin{proposition}[Finite \'etale schemes over a separably closed field]
\label{prop:av-finite-etale-points}
\lean{AlgebraicGeometry.Scheme.Hom.finrank_eq_natCard, AlgebraicGeometry.pointEquiv,
AlgebraicGeometry.natCard_pointHom_eq_finrank}
\leanok
Let $K$ be separably closed and let $X\to\Spec K$ be finite and \'etale.  Then
$X$ is affine, $\Gamma(X,\O_X)\cong K^{\pi_0}$ as a $K$-algebra for a finite
set $\pi_0$ in bijection with the points of $X$, and hence
$\deg(X\to\Spec K)=\#X$.  If moreover $K$ is algebraically closed, the
$K$-points of $X$ are exactly its points, so the number of $K$-points of $X$
equals $\deg(X\to\Spec K)$.
\end{proposition}
\begin{proof}
\leanok
A finite morphism to an affine scheme has affine source, so $X=\Spec S$ with
$S$ a finite \'etale $K$-algebra.  A finite \'etale algebra over a separably
closed field is a finite product of copies of the field: this is
\texttt{Algebra.FormallyEtale.equivPiOfIsSepClosed} in Mathlib.  The rank over
$K$ of $K^{\pi_0}$ is $\#\pi_0$, and the points of $\Spec K^{\pi_0}$ are the
$\#\pi_0$ primitive idempotents.  Over an algebraically closed field, a point
of a $K$-scheme locally of finite type whose residue field is $K$ is the same
thing as a $K$-point, which identifies $X(K)$ with $X$.
\end{proof}

\begin{theorem}[Recognising $(\Z/n)^r$]\label{thm:av-finite-abelian-group}
\lean{CommGroup.nonempty_equiv_pi_multiplicative_zmod_of_card_torsion,
AddCommGroup.nonempty_equiv_pi_zmod_of_card_torsion}
\leanok
Let $H$ be a finite abelian group, let $n\geq1$ and $r\geq0$, suppose that
$x^n=1$ for every $x\in H$, and suppose that for every $m$ dividing $n$ the
subgroup $\{x\in H:x^m=1\}$ has exactly $m^r$ elements.  Then $H\cong(\Z/n)^r$.
\end{theorem}
\begin{proof}
\leanok
By the structure theorem for finite abelian groups, $H\cong\prod_i\Z/p_i^{e_i}$
for primes $p_i$ and exponents $e_i\geq1$.  Each factor is killed by $n$, so
$p_i^{e_i}\mid n$.  Counting $m$-torsion in a product gives
$\#\{x\in H:x^m=1\}=\prod_i\gcd(m,p_i^{e_i})$.  Fix a prime $p$, let $k$ be
maximal with $p^k\mid n$ and take $m=p^k$; comparing with $m^r=p^{kr}$ forces
the number of indices $i$ with $p_i=p$ to be exactly $r$ and each of those
factors to be $\Z/p^k$.  Doing this for every prime dividing $n$ and
reassembling by the Chinese remainder theorem gives $H\cong(\Z/n)^r$.  This is
in the project's mirror file
\texttt{WeilConjectures/Mathlib/GroupTheory/FiniteAbelian/Basic.lean}.
\end{proof}

\section{The three propositions}\label{sec:av-three-propositions}

\begin{proposition}[Proposition 2.7: $[n]$ is an isogeny of degree $n^{2g}$]
\label{prop:av-isogeny}
\uses{def:av-mulby, def:av-degree, prop:av-basic}
\lean{AlgebraicGeometry.IsAbelianVariety.isFinite_mulBy,
AlgebraicGeometry.IsAbelianVariety.flat_mulBy,
AlgebraicGeometry.IsAbelianVariety.surjective_mulBy,
AlgebraicGeometry.IsAbelianVariety.finrank_mulBy}
Let $A$ be an abelian variety over a field $K$, smooth of relative dimension
$g$ over $K$, and let $n\neq0$ be a natural number.  Then $[n]\colon A\to A$ is
finite, flat and surjective, and $\deg[n]=n^{2g}$.
\end{proposition}
\begin{proof}
\uses{thm:av-mulby-etale-corollaries, thm:av-degree-formula,
thm:av-isogeny-general}
For $n$ invertible in $K$ the finiteness, flatness and surjectivity are proved
in \Cref{thm:av-mulby-etale-corollaries} by the deformation-theoretic route of
\Cref{sec:av-etale-route}, which uses no theory of line bundles.  The degree,
and the case where $n$ is divisible by the characteristic of $K$, are
\Cref{thm:av-degree-formula} and \Cref{thm:av-isogeny-general}; those rest on
the theorem of the cube and on intersection theory, and their five inputs are
listed as open obligations in \Cref{sec:av-cube}.
\end{proof}

\begin{proposition}[Proposition 2.8: the $n$-torsion is finite \'etale of rank
$n^{2g}$]
\label{prop:av-etale-torsion}
\uses{def:av-torsion, def:av-degree}
\lean{AlgebraicGeometry.IsAbelianVariety.isFinite_torsion,
AlgebraicGeometry.IsAbelianVariety.flat_torsion,
AlgebraicGeometry.IsAbelianVariety.finrank_torsion,
AlgebraicGeometry.IsAbelianVariety.etale_torsion,
AlgebraicGeometry.IsAbelianVariety.natCard_torsionHom,
AlgebraicGeometry.IsAbelianVariety.nonempty_torsionHom_mulEquiv_pi}
\leanok
Let $A$ be an abelian variety over $K$, smooth of relative dimension $g$ over
$K$, and let $n\neq0$ be invertible in $K$.  Then $A[n]\to\Spec K$ is finite,
flat and \'etale of rank $n^{2g}$.  If moreover $K$ is algebraically closed,
then $A[n](K)$ has exactly $n^{2g}$ elements and
\[
  A[n](K)\;\cong\;(\Z/n)^{2g}.
\]
\end{proposition}
\begin{proof}
\uses{prop:av-isogeny, lem:av-torsion-pullback, lem:av-degree-base-change,
thm:av-mulby-etale, prop:av-finite-etale-points, thm:av-finite-abelian-group}
\leanok
By \Cref{lem:av-torsion-pullback} the morphism $A[n]\to\Spec K$ is the base
change of $[n]$ along $e\colon\Spec K\to A$.  By \Cref{prop:av-isogeny} the
morphism $[n]$ is finite and flat of degree $n^{2g}$, and by
\Cref{thm:av-mulby-etale} it is \'etale because $n$ is invertible in $K$.  All
four properties are stable under base change, and the degree is preserved by
\Cref{lem:av-degree-base-change}; hence $A[n]\to\Spec K$ is finite, flat and
\'etale of rank $n^{2g}$.

Now let $K$ be algebraically closed.  By \Cref{prop:av-finite-etale-points}
applied to the finite \'etale $K$-scheme $A[n]$, the group $A[n](K)$ is finite
of cardinality equal to the rank of $A[n]$, that is $n^{2g}$.  By
\Cref{lem:av-torsion-pullback}, $A[n](K)=\{x\in A(K):x^n=1\}$; in particular
every element of $A[n](K)$ is killed by $n$, and for $m$ dividing $n$
\[
  \{x\in A[n](K):x^m=1\}=\{x\in A(K):x^m=1\}=A[m](K),
\]
because $x^m=1$ already implies $x^n=1$.  Such an $m$ is again invertible in
$K$, since a divisor of an integer invertible in a field is invertible in that
field, so the count just made applies to $m$ as well and gives
$\#A[m](K)=m^{2g}$.  Thus $A[n](K)$ is a finite abelian group killed by $n$
whose $m$-torsion has exactly $m^{2g}$ elements for every $m$ dividing $n$, and
\Cref{thm:av-finite-abelian-group} gives $A[n](K)\cong(\Z/n)^{2g}$.
\end{proof}

\begin{proposition}[Proposition 2.9: $A(K)$ is $n$-divisible]
\label{prop:av-divisibility}
\uses{def:av-abelian-variety}
\lean{AlgebraicGeometry.IsAbelianVariety.exists_npow_eq,
AlgebraicGeometry.IsAbelianVariety.rootableBy}
\leanok
Let $A$ be an abelian variety over an algebraically closed field $K$ and let
$n\neq0$.  Then every element of $A(K)$ is an $n$-th power; equivalently
$A(K)/n\cdot A(K)=0$.
\end{proposition}
\begin{proof}
\uses{prop:av-isogeny, lem:av-mulby-points, prop:av-basic}
\leanok
By \Cref{prop:av-isogeny} the morphism $[n]\colon A\to A$ is finite and
surjective.  Let $x\in A(K)$.  Since $K$ is algebraically closed and $A$ is
locally of finite type over $K$, the $K$-points of $A$ are exactly the closed
points of the scheme $A$; write $x_0\in A$ for the closed point corresponding
to $x$.  By surjectivity there is $z\in A$ with $[n](z)=x_0$.

The point $z$ is closed.  Indeed $\{x_0\}$ is closed, so $[n]^{-1}(\{x_0\})$ is
closed, and it is finite because $[n]$ is a finite morphism; and $A$ is a
Jacobson space, being locally of finite type over a field, so a finite closed
subset of $A$ consists of closed points.  Hence $z$ is a closed point of $A$
and therefore corresponds to a $K$-point $y\in A(K)$ with $[n]\circ y=x$.  By
\Cref{lem:av-mulby-points}, $[n]\circ y=y^n$, so $y^n=x$.
\end{proof}

\begin{remark}[What the \'etale route does and does not give]
\label{rem:av-etale-route-scope}
\uses{prop:av-isogeny, thm:av-mulby-etale-corollaries}
\Cref{sec:av-etale-route} proves finiteness, flatness and surjectivity of $[n]$
for $n$ invertible in $K$, over an arbitrary field.  It does \emph{not} prove
that $\deg[n]=n^{2g}$: nothing in a deformation-theoretic argument sees the
dimension $g$.  Consequently the rank statement of
\Cref{prop:av-etale-torsion}, and with it the isomorphism
$A[n](K)\cong(\Z/n)^{2g}$, genuinely depends on \Cref{sec:av-cube}, whereas the
finiteness and \'etaleness of $A[n]$ do not.
\end{remark}

\section{The \'etale route}\label{sec:av-etale-route}

Throughout this section $A$ is an abelian variety over a field $K$ and $n$ is a
natural number that is invertible in $K$, that is, $n\neq0$ in $K$.  No
hypothesis is placed on $K$ unless it is stated; \Cref{rem:av-field-hypotheses}
records, statement by statement, which of them need $K$ algebraically closed,
which need $K$ perfect and which hold over an arbitrary field.  It will turn out
that everything from \Cref{sec:av-etale-kernel} onwards holds over an arbitrary
field.

The plan is: the kernel of $A(T)\to A(T_0)$ along a square-zero thickening is a
$K$-vector space (\Cref{sec:av-etale-kernel}); therefore $x\mapsto x^n$ is
bijective on it; therefore $[n]$ is formally unramified and formally smooth
(\Cref{sec:av-etale-unramified,sec:av-etale-smooth}); therefore $[n]$ is
\'etale (\Cref{sec:av-etale-conclusion}); therefore, using properness of $A$
and Zariski's main theorem, $[n]$ is finite, flat and surjective
(\Cref{sec:av-etale-corollaries}).

\statusnote{The results of this section are formalised, but not by the
decomposition written here.  The Lean sources are
\texttt{WeilConjectures/AlgebraicGeometry/Group/Infinitesimal.lean}, which
proves \Cref{cor:av-kernel-npow} for an arbitrary group scheme over $K$,
commutative or not, and
\texttt{WeilConjectures/AlgebraicGeometry/Group/Etale.lean}, which proves
\Cref{thm:av-mulby-formally-unramified} for a commutative group scheme over
$K$ and \Cref{thm:av-mulby-formally-smooth} for a commutative group scheme
that is moreover smooth over $K$.  None of the three needs properness or
integrality, so all of them hold well beyond abelian varieties; restating them
here in that generality is left for later.  There are three points where the
Lean route departs from the text below.  \emph{First}, \Cref{cor:av-kernel-npow}
is proved directly, by a product formula for ring maps out of an affine chart
around the identity, with no cotangent space and no module structure anywhere;
so the statement of \Cref{cor:av-kernel-npow} is formalised (as
\texttt{AlgebraicGeometry.GrpObj.eq\_one\_of\_pow\_eq\_one} and
\texttt{AlgebraicGeometry.GrpObj.exists\_pow\_eq}) while the cotangent-space
material of \Cref{sec:av-etale-kernel} that the proof below runs through ---
\Cref{def:av-cotangent} to \Cref{cor:av-kernel-module} --- has no Lean
counterpart and is not on the path to \Cref{thm:av-mulby-etale}.  That direct
argument is not written out in this chapter; writing it out, and keeping the
cotangent route as the classical account, is a follow-up.  \emph{Second},
\Cref{thm:av-mulby-formally-smooth} and \Cref{lem:av-mulby-lfp} are packaged
into the single declaration
\texttt{AlgebraicGeometry.GrpObj.smooth\_npow\_left}, because Mathlib's
\texttt{AlgebraicGeometry.Smooth} is formal smoothness together with local
finite presentation.  That packaging also saddles \Cref{lem:av-mulby-lfp} with
the hypothesis that $n$ be invertible in $K$, which it does not need; an
unconditional Lean lemma for it is a follow-up.  \emph{Third}, the three open
obligations formerly recorded at \Cref{thm:av-mulby-formally-unramified},
\Cref{thm:av-mulby-formally-smooth} and \Cref{lem:av-mulby-lfp} are removed,
but the Mathlib gaps they recorded are circumvented rather than filled:
Mathlib \texttt{v4.33.1} still has no infinitesimal lifting criterion for
\texttt{AlgebraicGeometry.FormallyUnramified}, no
\texttt{AlgebraicGeometry.FormallySmooth} and no cancellation instance for
\texttt{LocallyOfFinitePresentation}.  The formalisation sidesteps the first
two by working on affine charts, in the form described in
\Cref{rem:av-unramified-affine-local} and \Cref{rem:av-smooth-affine-local},
and the third by cancelling finite presentation at ring level with
\texttt{RingHom.FinitePresentation.of\_comp\_finiteType}.}

\subsection{Translations}\label{sec:av-etale-translations}

\begin{construction}[Translation by a rational point]\label{def:av-translation}
\uses{def:av-abelian-variety}
Let $G$ be a group scheme over $K$ with structure morphism $\pi\colon G\to\Spec
K$, and let $a\in G(K)$.  Write $a_G:=a\circ\pi\in G(G)$ for the constant
$G$-point at $a$.  The \emph{translation} $t_a\colon G\to G$ is the $G$-point
$a_G\cdot\mathrm{id}$ of $G$; its effect on points is $x\mapsto a\cdot x$.
\end{construction}

\begin{lemma}[Translations are isomorphisms]\label{lem:av-translation}
\uses{def:av-translation}
Let $G$ be a group scheme over $K$ and $a,b\in G(K)$.  Then $t_a\circ
t_b=t_{ab}$ and $t_1=\mathrm{id}$; in particular $t_a$ is an isomorphism with
inverse $t_{a^{-1}}$.  This holds over an arbitrary field.
\end{lemma}
\begin{proof}
For a $K$-scheme $T$ write $a_T:=a\circ(T\to\Spec K)\in G(T)$; then
$a\mapsto a_T$ is a group homomorphism $G(K)\to G(T)$.  Directly from the
construction, $t_a\circ u=a_T\cdot u$ for every $u\in G(T)$.  Hence for every
$T$ and every $u\in G(T)$,
\[
  (t_a\circ t_b)\circ u=t_a\circ(t_b\circ u)=a_T\cdot(b_T\cdot u)
  =(ab)_T\cdot u=t_{ab}\circ u .
\]
Taking $T=G$ and $u=\mathrm{id}$ gives $t_a\circ t_b=t_{ab}$.  Also
$t_1\circ u=1\cdot u=u$, so $t_1=\mathrm{id}$.  Taking $b=a^{-1}$ gives the
last claim.
\end{proof}

\begin{lemma}[Translations intertwine $[n]$]
\label{lem:av-translation-intertwines}
\uses{def:av-translation, def:av-mulby}
Let $G$ be a \emph{commutative} group scheme over $K$, let $a\in G(K)$ and let
$n\in\N$.  Then
\[
  [n]\circ t_a\;=\;t_{a^n}\circ[n].
\]
This holds over an arbitrary field.
\end{lemma}
\begin{proof}
\uses{lem:av-mulby-points, lem:av-translation}
For $u\in G(T)$ one has, using \Cref{lem:av-mulby-points} and the notation of
the previous proof,
\[
  ([n]\circ t_a)\circ u=(a_T\cdot u)^n=(a_T)^n\cdot u^n=(a^n)_T\cdot u^n
  =(t_{a^n}\circ[n])\circ u,
\]
the middle equality because $G(T)$ is commutative.  Take $T=G$ and
$u=\mathrm{id}$.
\end{proof}

\begin{lemma}[Spreading a property of $[n]$ out by translation]
\label{lem:av-translation-spreading}
\uses{lem:av-translation-intertwines}
Let $A$ be an abelian variety over $K$ and let $P$ be a property of morphisms
of schemes that is local on source and target and stable under pre- and
post-composition with isomorphisms.  Let
\[
  V:=\{x\in A:\ [n]\ \text{has}\ P\ \text{at}\ x\}.
\]
Then $t_a(V)=V$ for every $a\in A(K)$.  If moreover $V$ is open and non-empty
and $K$ is algebraically closed, then $V=A$.
\end{lemma}
\begin{proof}
\uses{lem:av-translation, lem:av-translation-intertwines}
\emph{Stability under translation.}  By
\Cref{lem:av-translation-intertwines}, $[n]=t_{a^n}\circ[n]\circ t_a^{-1}$, and
$t_a$ and $t_{a^n}$ are isomorphisms by \Cref{lem:av-translation}.  Since $P$
is local and stable under composing with isomorphisms on either side, $[n]$ has
$P$ at $t_a(x)$ if and only if it has $P$ at $x$; that is $t_a(V)=V$.

\emph{Conclusion when $K=\bar K$.}  Over an algebraically closed field the
closed points of $A$ are exactly the $K$-points.  Pick a $K$-point $v$ whose
image lies in $V$ --- possible because $V$ is a non-empty open of a scheme
locally of finite type over $K$, hence contains a closed point.  Let $x$ be any
closed point of $A$, corresponding to $b\in A(K)$.  Then
$x=t_{bv^{-1}}(v)\in t_{bv^{-1}}(V)=V$.  So $V$ contains every closed point of
$A$; its complement is closed and contains no closed point, hence is empty
because $A$ is a Jacobson space.  Thus $V=A$.
\end{proof}

\begin{remark}[Translations are not used below]\label{rem:av-translation-unused}
\uses{lem:av-translation-spreading}
\Cref{lem:av-translation-spreading} is the classical device for reducing a
statement about $[n]$ to its behaviour at the identity, and it is what one uses
if one computes the differential of $[n]$ at the identity and wants to conclude
everywhere.  The route taken in
\Cref{sec:av-etale-kernel,sec:av-etale-unramified,sec:av-etale-smooth,sec:av-etale-conclusion}
does not use it: those arguments treat all points at once, being statements
about lifting $T$-points for arbitrary affine $T$, and they need no hypothesis
on $K$.  This matters for the formalisation, because
\Cref{lem:av-translation-spreading} is the only step of this section that would
require $K$ to be algebraically closed, while the Lean statements are over an
arbitrary field.  The translation lemmas are recorded here because they are
cheap, because they are needed to compare with the classical accounts, and
because \Cref{lem:av-translation-intertwines} is used in \Cref{sec:av-cube} to
see that every fibre of $[n]$ is a translate of $A[n]$.
\end{remark}

\subsection{The kernel of $A(R)\to A(R/I)$}\label{sec:av-etale-kernel}

This subsection is the one step of the \'etale route that is real work.  Fix a
group scheme $G$ over $K$; properness, integrality and smoothness are \emph{not}
used anywhere in this subsection.  We describe the kernel of the restriction map
on points along a square-zero thickening, and in particular give it a
$K$-vector space structure.

\begin{definition}[Cotangent space at a rational point]\label{def:av-cotangent}
Let $X$ be a $K$-scheme and $x\in X(K)$ a $K$-rational point, with image the
point $\xi\in X$; write $\m_\xi\subseteq\O_{X,\xi}$ for the maximal ideal, so
that $\O_{X,\xi}/\m_\xi=K$.  The \emph{cotangent space} of $X$ at $x$ is the
$K$-vector space
\[
  \mathfrak c(X,x):=\m_\xi/\m_\xi^2 .
\]
A morphism $f\colon X\to Y$ of $K$-schemes with $f\circ x=y$ induces a
$K$-linear map $\mathfrak c(f)\colon\mathfrak c(Y,y)\to\mathfrak c(X,x)$,
because $f$ induces a local homomorphism $\O_{Y,\eta}\to\O_{X,\xi}$ of local
rings.  Thus $\mathfrak c(-)$ is a contravariant functor on $K$-schemes
equipped with a $K$-rational point.
\end{definition}

\begin{lemma}[Cotangent space computed on an affine open]
\label{lem:av-cotangent-affine}
\uses{def:av-cotangent}
In the situation of \Cref{def:av-cotangent}, let $U=\Spec S\subseteq X$ be an
affine open containing $\xi$ and let $\m:=\ker(x^\sharp\colon S\to K)$ be the
corresponding maximal ideal.  Then the localisation map induces an isomorphism
of $K$-vector spaces
\[
  \m/\m^2\;\xrightarrow{\ \sim\ }\;\mathfrak c(X,x).
\]
In particular $\mathfrak c(X,x)$ may be computed in any affine open containing
$\xi$.
\end{lemma}
\begin{proof}
Localisation is exact, so $(\m/\m^2)_\m=\m_\xi/\m_\xi^2$.  The $S$-module
$M:=\m/\m^2$ is annihilated by $\m$, hence is a module over $S/\m=K$; and an
element $s\in S\setminus\m$ acts on $M$ as multiplication by its class in
$S/\m=K$, which is non-zero, hence invertibly.  Therefore $M\to M_\m$ is an
isomorphism.
\end{proof}

\begin{lemma}[Cotangent space of a product]\label{lem:av-cotangent-product}
\uses{def:av-cotangent}
Let $X,Y$ be $K$-schemes with $K$-points $x,y$, and let $(x,y)$ be the induced
$K$-point of $X\times_KY$.  Then the two projections $p_1,p_2$ induce an
isomorphism
\[
  \bigl(\mathfrak c(p_1),\mathfrak c(p_2)\bigr)\colon
  \mathfrak c(X,x)\oplus\mathfrak c(Y,y)\;\xrightarrow{\ \sim\ }\;
  \mathfrak c(X\times_KY,(x,y)),
\]
whose inverse has components $\mathfrak c(j_1)$ and $\mathfrak c(j_2)$, where
$j_1:=(\mathrm{id}_X,y)\colon X\to X\times_KY$ and
$j_2:=(x,\mathrm{id}_Y)\colon Y\to X\times_KY$.
\end{lemma}
\begin{proof}
\uses{lem:av-cotangent-affine}
By \Cref{lem:av-cotangent-affine} we may compute in affine opens: choose
$U=\Spec S$ containing the image of $x$ and $V=\Spec S'$ containing the image
of $y$, so that $U\times_KV=\Spec(S\otimes_KS')$ is an affine open of
$X\times_KY$ containing the image of $(x,y)$.  Write $\varepsilon\colon S\to K$
and $\varepsilon'\colon S'\to K$ for the two augmentations,
$\m=\ker\varepsilon$ and $\m'=\ker\varepsilon'$.  As $K$-modules $S=K\oplus\m$
and $S'=K\oplus\m'$, so
\[
  S\otimes_KS'=K\ \oplus\ (\m\otimes_KK)\ \oplus\ (K\otimes_K\m')\ \oplus\
  (\m\otimes_K\m'),
\]
and the augmentation ideal $\mathfrak n:=\ker(\varepsilon\otimes\varepsilon')$
is the sum of the last three summands.  Multiplying out and using
$(\m\otimes1)\cdot(1\otimes\m')=\m\otimes\m'$,
\[
  \mathfrak n^2=(\m^2\otimes_KK)\ \oplus\ (\m\otimes_K\m')\ \oplus\
  (K\otimes_K(\m')^2).
\]
Hence $\mathfrak n/\mathfrak n^2\cong\m/\m^2\oplus\m'/(\m')^2$.  The ring maps of
$p_1$ and $p_2$ are $s\mapsto s\otimes1$ and $t\mapsto1\otimes t$, so under
this decomposition $\mathfrak c(p_1)$ and $\mathfrak c(p_2)$ are the two
inclusions; the ring maps of $j_1$ and $j_2$ are
$s\otimes t\mapsto\varepsilon'(t)s$ and $s\otimes t\mapsto\varepsilon(s)t$, so
$\mathfrak c(j_1)$ and $\mathfrak c(j_2)$ are the two projections.  These are
mutually inverse.
\end{proof}

\begin{lemma}[The cotangent space sees the group law]
\label{lem:av-cotangent-group-law}
\uses{lem:av-cotangent-product, def:av-abelian-variety}
Let $G$ be a group scheme over $K$ with identity section $e$, and write
$\mathfrak c:=\mathfrak c(G,e)$.  Identify $\mathfrak c(G\times_KG,(e,e))$ with
$\mathfrak c\oplus\mathfrak c$ by \Cref{lem:av-cotangent-product}.  Then the
multiplication $\mu$ induces
\[
  \mathfrak c(\mu)\colon\mathfrak c\longrightarrow\mathfrak c\oplus\mathfrak c,
  \qquad \mathfrak c(\mu)(s)=(s,s).
\]
\end{lemma}
\begin{proof}
\uses{lem:av-cotangent-product}
The morphism $\mu$ sends the $K$-point $(e,e)$ to $e$, so $\mathfrak c(\mu)$ is
defined.  By \Cref{lem:av-cotangent-product} the components of an element of
$\mathfrak c(G\times_KG,(e,e))$ are read off by applying $\mathfrak c(j_1)$ and
$\mathfrak c(j_2)$, where $j_1=(\mathrm{id}_G,e)$ and $j_2=(e,\mathrm{id}_G)$.
So the two components of $\mathfrak c(\mu)$ are $\mathfrak c(j_1)\circ\mathfrak
c(\mu)=\mathfrak c(\mu\circ j_1)$ and likewise $\mathfrak c(\mu\circ j_2)$.
The unit axioms of a group object say exactly $\mu\circ
j_1=\mathrm{id}_G=\mu\circ j_2$, so by functoriality both components are the
identity of $\mathfrak c$.
\end{proof}

\begin{lemma}[The kernel along a square-zero thickening]
\label{lem:av-kernel-cotangent}
\uses{def:av-cotangent, lem:av-cotangent-affine}
Let $G$ be a group scheme over $K$, let $R$ be a $K$-algebra and let
$I\subseteq R$ be an ideal with $I^2=0$; put $R_0:=R/I$, $T:=\Spec R$ and
$T_0:=\Spec R_0$.  Write
\[
  N(R,I):=\ker\bigl(G(R)\longrightarrow G(R_0)\bigr).
\]
Choose an affine open $U=\Spec S$ of $G$ containing the image $\eta$ of $e$,
let $\varepsilon\colon S\to K$ be the augmentation given by $e$, let
$\m=\ker\varepsilon$, and let $\varepsilon_R\colon S\to R$ be $\varepsilon$
followed by the structure map $K\to R$.  Then every $w\in N(R,I)$ factors
through $U$ and so is given by a $K$-algebra map $\varphi_w\colon S\to R$; the
difference $D_w:=\varphi_w-\varepsilon_R$ takes values in $I$, vanishes on
$K\subseteq S$ and on $\m^2$, and the assignment
\[
  \Theta\colon w\longmapsto\bigl(\text{the induced map }
  \m/\m^2=\mathfrak c(G,e)\to I\bigr)
\]
is a bijection
$\Theta\colon N(R,I)\xrightarrow{\ \sim\ }\Hom_K(\mathfrak c(G,e),I)$.
\end{lemma}
\begin{proof}
\uses{lem:av-cotangent-affine}
\emph{Step 1: $w$ factors through $U$.}  Since $I$ is nilpotent, the closed
immersion $T_0\hookrightarrow T$ is a homeomorphism on underlying spaces.  The
restriction of $w$ to $T_0$ is the trivial point, that is $e$ composed with
$T_0\to\Spec K$, whose image is $\{\eta\}$.  Hence the image of the continuous
map underlying $w$ is $\{\eta\}\subseteq U$, so $w^{-1}(U)=T$ and $w$ factors
uniquely through the open subscheme $U$.  As $U$ is affine, morphisms $T\to U$
correspond bijectively to $K$-algebra maps $S\to R$; write $\varphi_w$ for the
one corresponding to $w$.

\emph{Step 2: $D_w$ is a derivation into $I$.}  That $w$ restricts to the
trivial point on $T_0$ says precisely that $\varphi_w$ followed by $R\to R_0$
equals $\varepsilon$ followed by $K\to R_0$; that is,
$D_w:=\varphi_w-\varepsilon_R$ takes values in $I$.  It is $K$-linear, and for
$s,t\in S$
\[
  D_w(st)=\varphi_w(s)\varphi_w(t)-\varepsilon_R(s)\varepsilon_R(t)
  =\varepsilon_R(s)D_w(t)+\varepsilon_R(t)D_w(s)+D_w(s)D_w(t),
\]
where the last term lies in $I^2=0$.  So $D_w$ is a $K$-derivation of $S$ into
$I$, the $S$-module structure on $I$ being through $\varepsilon$.  In
particular $D_w(1)=\varphi_w(1)-1=0$, so $D_w$ vanishes on the image of $K$,
and $D_w(\m^2)=0$ because $\varepsilon(\m)=0$.  Since $S=K\oplus\m$ as a
$K$-module, $D_w$ is determined by the induced $K$-linear map $\m/\m^2\to I$,
which is $\Theta(w)$; it is a map out of $\mathfrak c(G,e)$ by
\Cref{lem:av-cotangent-affine}.

\emph{Step 3: $\Theta$ is bijective.}  Injectivity is Step 2: $D_w$, hence
$\varphi_w$, hence $w$, is determined by $\Theta(w)$.  For surjectivity let
$\delta\colon\m/\m^2\to I$ be $K$-linear, define $\tilde\delta\colon S\to I$ as
the composite
$S=K\oplus\m\twoheadrightarrow\m\twoheadrightarrow\m/\m^2
\xrightarrow{\ \delta\ }I$
and put $\varphi:=\varepsilon_R+\tilde\delta$.  Then $\varphi$ is $K$-linear
and $\varphi(1)=1$.  For $s,t\in S$ write $s=\varepsilon(s)+s_0$ and
$t=\varepsilon(t)+t_0$ with $s_0,t_0\in\m$; then
$st=\varepsilon(s)\varepsilon(t)+\varepsilon(s)t_0+\varepsilon(t)s_0+s_0t_0$
with $s_0t_0\in\m^2$, so
\[
  \tilde\delta(st)=\varepsilon(s)\tilde\delta(t)+\varepsilon(t)\tilde\delta(s),
\]
and therefore
\[
  \varphi(s)\varphi(t)
  =\varepsilon_R(s)\varepsilon_R(t)+\varepsilon_R(s)\tilde\delta(t)
   +\varepsilon_R(t)\tilde\delta(s)+\tilde\delta(s)\tilde\delta(t)
  =\varepsilon_R(st)+\tilde\delta(st)=\varphi(st),
\]
the term $\tilde\delta(s)\tilde\delta(t)$ vanishing because $I^2=0$.  So
$\varphi$ is a $K$-algebra map; the corresponding morphism $T\to U\subseteq G$
restricts on $T_0$ to the trivial point because $\tilde\delta$ has values in
$I$, hence lies in $N(R,I)$, and $\Theta$ sends it to $\delta$.
\end{proof}

\begin{lemma}[$\Theta$ is a group isomorphism]\label{lem:av-kernel-group-law}
\uses{lem:av-kernel-cotangent, lem:av-cotangent-group-law}
In the situation of \Cref{lem:av-kernel-cotangent}, $N(R,I)$ is a subgroup of
$G(R)$ and
\[
  \Theta\colon N(R,I)\;\xrightarrow{\ \sim\ }\;
  \bigl(\Hom_K(\mathfrak c(G,e),I),\,+\bigr)
\]
is an isomorphism of groups.  In particular $N(R,I)$ is abelian, whether or not
$G$ is.
\end{lemma}
\begin{proof}
\uses{lem:av-kernel-cotangent, lem:av-cotangent-group-law,
lem:av-cotangent-product}
\emph{Step 0: naturality of $\Theta$.}  Let $h\colon G'\to G$ be a morphism of
$K$-schemes with $h\circ e'=e$, where $e'$ is the identity section of a group
scheme $G'$.  If $w'\in G'(R)$ restricts on $T_0$ to the trivial point, then so
does $h\circ w'$, and directly from the description of $\Theta$ as ``the map
induced by $\varphi-\varepsilon_R$ on $\m/\m^2$'' one gets
\[
  \Theta(h\circ w')=\Theta(w')\circ\mathfrak c(h),
\]
where $\mathfrak c(h)\colon\mathfrak c(G,e)\to\mathfrak c(G',e')$.  Indeed the
ring map of $h\circ w'$ is $\varphi_{w'}\circ h^\sharp$, so its difference from
the augmentation is $D_{w'}\circ h^\sharp$, and passing to $\m/\m^2$ turns
$h^\sharp$ into $\mathfrak c(h)$.  Here $\Theta$ on the left is the map of
\Cref{lem:av-kernel-cotangent} for $G$ and on the right the one for $G'$; note
that \Cref{lem:av-kernel-cotangent} used nothing about $G$ beyond the pointed
scheme $(G,e)$, so it applies to $G\times_KG$ with the point $(e,e)$ as well.

\emph{Step 1: the pair $(w_1,w_2)$.}  That $N(R,I)$ is a subgroup is clear,
being the kernel of the group homomorphism $G(R)\to G(R_0)$.  Let
$w_1,w_2\in N(R,I)$ and put $D_i:=\Theta(w_i)$.  The morphism
$(w_1,w_2)\colon T\to G\times_KG$ restricts on $T_0$ to the trivial point of
$G\times_KG$, so it lies in the corresponding kernel and has an invariant
\[
  \beta:=\Theta\bigl((w_1,w_2)\bigr)\in
  \Hom_K\bigl(\mathfrak c(G\times_KG,(e,e)),I\bigr).
\]
Applying Step 0 to $h=p_1$ and $h=p_2$ gives
$\beta\circ\mathfrak c(p_1)=\Theta(p_1\circ(w_1,w_2))=\Theta(w_1)=D_1$ and
likewise $\beta\circ\mathfrak c(p_2)=D_2$.  Under the identification
$\mathfrak c(G\times_KG,(e,e))\cong\mathfrak c\oplus\mathfrak c$ of
\Cref{lem:av-cotangent-product}, whose two inclusions are $\mathfrak c(p_1)$
and $\mathfrak c(p_2)$, this says $\beta(a,b)=D_1(a)+D_2(b)$.

\emph{Step 2: conclusion.}  By definition $w_1\cdot w_2=\mu\circ(w_1,w_2)$, so
Step 0 applied to $h=\mu$ and \Cref{lem:av-cotangent-group-law} give
\[
  \Theta(w_1\cdot w_2)=\beta\circ\mathfrak c(\mu)
  =\bigl(s\mapsto\beta(s,s)\bigr)=D_1+D_2=\Theta(w_1)+\Theta(w_2).
\]
So $\Theta$ is a group homomorphism, and it is bijective by
\Cref{lem:av-kernel-cotangent}.  A group isomorphic to an abelian group is
abelian.
\end{proof}

\begin{corollary}[The kernel is an $R_0$-module, in particular a $K$-vector
space]
\label{cor:av-kernel-module}
\uses{lem:av-kernel-group-law}
In the situation of \Cref{lem:av-kernel-cotangent}, $N(R,I)$ carries a natural
structure of module over $R_0=R/I$, in particular of $K$-vector space, whose
underlying additive group is $N(R,I)$ written additively.
\end{corollary}
\begin{proof}
\uses{lem:av-kernel-group-law}
The ideal $I$ is killed by $I$, so its $R$-module structure factors through
$R_0$.  The $R_0$-action on $I$ is $K$-linear, so it induces an $R_0$-module
structure on $\Hom_K(\mathfrak c(G,e),I)$ by acting on the target, compatible
with the addition used in \Cref{lem:av-kernel-group-law}.  Transport this
structure along the group isomorphism $\Theta$.  The $K$-vector space structure
is its restriction along $K\to R_0$; explicitly, $\lambda\cdot w$ is the
element of $N(R,I)$ with $\Theta(\lambda\cdot w)=\lambda\,\Theta(w)$.
\end{proof}

\begin{corollary}[$n$-th powers on the kernel]\label{cor:av-kernel-npow}
\lean{AlgebraicGeometry.GrpObj.eq_one_of_pow_eq_one,
AlgebraicGeometry.GrpObj.exists_pow_eq}
\leanok
Let $G$ be a group scheme over $K$, let $R$ be a $K$-algebra, let $I\subseteq R$
be an ideal with $I^2=0$, let $n$ be invertible in $K$, and write
$N(R,I)=\ker\bigl(G(R)\to G(R/I)\bigr)$.  Then
\begin{enumerate}
  \item if $w\in N(R,I)$ satisfies $w^n=1$ then $w=1$; and
  \item every $w\in N(R,I)$ is the $n$-th power of an element of $N(R,I)$.
\end{enumerate}
Since $N(R,I)$ is abelian (\Cref{lem:av-kernel-group-law}), the two together
say that
\[
  N(R,I)\longrightarrow N(R,I),\qquad w\longmapsto w^n,
\]
is bijective.  This holds over an arbitrary field.
\end{corollary}
\begin{proof}
\uses{cor:av-kernel-module, lem:av-kernel-group-law}
Under the group isomorphism $\Theta$ of \Cref{lem:av-kernel-group-law} the map
$w\mapsto w^n$ becomes $D\mapsto nD$ on $\Hom_K(\mathfrak c(G,e),I)$, and by
\Cref{cor:av-kernel-module} this is multiplication by the image of $n$ in $K$,
which is a unit.  So it is bijective, and in particular (1) and (2) hold.
\end{proof}
\statusnote{This is not the Lean proof.  In Lean (1) and (2) are proved
directly, in an affine chart around the identity and with no commutativity
assumption on $G$, so the cotangent space plays no part; see the status note at
the head of \Cref{sec:av-etale-route}.}

\begin{remark}[Where the vector space structure comes from]
\label{rem:av-kernel-source}
\uses{cor:av-kernel-module}
It is worth isolating the point, because it is the only place in this section
where anything happens.  The group $N(R,I)$ is a group of $R$-points; there is
no scalar multiplication on $A(R)$ and none on $A(R_0)$.  Scalars appear only
after $N(R,I)$ has been identified with
$\Hom_K(\mathfrak c(A,e),I)$, and there they come from the \emph{target} $I$:
an ideal with $I^2=0$ in a $K$-algebra is a module over $R/I$ and \emph{a
fortiori} a $K$-vector space, and $\Hom_K(-,I)$ inherits that.  The content of
\Cref{lem:av-kernel-group-law} is that the multiplication of $A(R)$, restricted
to $N(R,I)$, becomes the addition of $\Hom_K(-,I)$; and that in turn is the
unit axiom of the group object, read on cotangent spaces
(\Cref{lem:av-cotangent-group-law}).  Nothing about $A$ beyond its being a
group scheme over $K$ is used --- not properness, not integrality, not even
smoothness.
\end{remark}

\subsection{$[n]$ is formally unramified}\label{sec:av-etale-unramified}

\begin{lemma}[Lifts land in a prescribed affine open]
\label{lem:av-lift-lands-in-open}
Let $T_0\hookrightarrow T$ be a closed immersion of schemes defined by a
nilpotent ideal, let $X$ be a scheme, let $a\colon T_0\to X$ be a morphism
whose image lies in an open subscheme $U\subseteq X$, and let $u\colon T\to X$
satisfy $u|_{T_0}=a$.  Then $u$ factors through $U$.
\end{lemma}
\begin{proof}
The closed immersion $T_0\hookrightarrow T$ is a homeomorphism on underlying
spaces because its ideal is nilpotent.  Hence the image of the continuous map
underlying $u$ equals the image of the one underlying $a$, which lies in $U$.
So $u^{-1}(U)=T$ and $u$ factors uniquely through the open subscheme $U$.
\end{proof}

\begin{theorem}[$[n]$ is formally unramified]
\label{thm:av-mulby-formally-unramified}
\uses{def:av-mulby, cor:av-kernel-npow}
\lean{AlgebraicGeometry.GrpObj.formallyUnramified_npow_left}
\leanok
Let $A$ be an abelian variety over a field $K$ and let $n$ be invertible in
$K$.  Let $R$ be a $K$-algebra, let $I\subseteq R$ be an ideal with $I^2=0$,
and put $T=\Spec R$ and $T_0=\Spec R/I$.  Let $u,v\colon T\to A$ be morphisms
of $K$-schemes with
\[
  [n]\circ u=[n]\circ v\qquad\text{and}\qquad u|_{T_0}=v|_{T_0}.
\]
Then $u=v$.  Consequently $[n]$ is formally unramified.  This holds over an
arbitrary field.
\end{theorem}
\begin{proof}
\uses{cor:av-kernel-npow, lem:av-mulby-points, prop:av-basic,
lem:av-kernel-group-law}
Work in the group $A(R)=\Hom_K(T,A)$, which is commutative by
\Cref{prop:av-basic}.  Put $w:=u\cdot v^{-1}\in A(R)$.  Restriction to $T_0$ is
a group homomorphism $A(R)\to A(R/I)$, and $u|_{T_0}=v|_{T_0}$, so $w$ lies in
$N(R,I)=\ker\bigl(A(R)\to A(R/I)\bigr)$.  By \Cref{lem:av-mulby-points},
$[n]\circ u=u^n$ and $[n]\circ v=v^n$, so $u^n=v^n$, and by commutativity
\[
  w^n=(u\cdot v^{-1})^n=u^n\cdot(v^n)^{-1}=1 .
\]
By \Cref{cor:av-kernel-npow} the map $w\mapsto w^n$ is injective on $N(R,I)$
and sends $1$ to $1$, so $w=1$, that is $u=v$.

For the last sentence: formal unramifiedness of $[n]$ is the assertion that in
every commutative square consisting of a square-zero thickening
$T_0\hookrightarrow T$ of affine schemes, a morphism $T_0\to A$ into the source
copy and a morphism $T\to A$ into the target copy, there is at most one
diagonal filler; a diagonal filler is a morphism $u\colon T\to A$ with
$u|_{T_0}$ and $[n]\circ u$ prescribed, so two of them satisfy the hypotheses
above and therefore coincide.
\end{proof}

\begin{remark}[Reading this affine-locally]
\label{rem:av-unramified-affine-local}
\uses{thm:av-mulby-formally-unramified, lem:av-lift-lands-in-open}
Mathlib's \texttt{AlgebraicGeometry.FormallyUnramified} is obtained by
\texttt{HasRingHomProperty} from \texttt{RingHom.FormallyUnramified}, so
concretely one must show, for affine opens $U,V\subseteq A$ with
$[n](U)\subseteq V$, that $\Gamma(U,\O)$ is formally unramified over
$\Gamma(V,\O)$: for every $\Gamma(V,\O)$-algebra $B$ and every square-zero
ideal $I\subseteq B$, the restriction map
$\Hom_{\Gamma(V,\O)}(\Gamma(U,\O),B)\to\Hom_{\Gamma(V,\O)}(\Gamma(U,\O),B/I)$
is injective.  Two elements of the source are two morphisms
$u,v\colon\Spec B\to U$ that agree on $\Spec B/I$ and have the same composite
with $[n]$; since $U\hookrightarrow A$ is a monomorphism it suffices to prove
$u=v$ as morphisms to $A$, which is \Cref{thm:av-mulby-formally-unramified}.
\end{remark}

\subsection{$[n]$ is formally smooth}\label{sec:av-etale-smooth}

\begin{theorem}[$[n]$ is formally smooth]\label{thm:av-mulby-formally-smooth}
\uses{def:av-mulby, cor:av-kernel-npow, prop:av-basic}
\lean{AlgebraicGeometry.GrpObj.smooth_npow_left}
\leanok
Let $A$ be an abelian variety over a field $K$ and let $n$ be invertible in
$K$.  Let $R$ be a $K$-algebra, let $I\subseteq R$ be an ideal with $I^2=0$,
and put $T=\Spec R$ and $T_0=\Spec R/I$.  Let $a\colon T_0\to A$ and
$b\colon T\to A$ be morphisms of $K$-schemes with $[n]\circ a=b|_{T_0}$.  Then
there is $u\colon T\to A$ with
\[
  u|_{T_0}=a\qquad\text{and}\qquad[n]\circ u=b .
\]
This holds over an arbitrary field.
\end{theorem}
\begin{proof}
\uses{cor:av-kernel-npow, lem:av-mulby-points, prop:av-basic,
lem:av-lift-lands-in-open, lem:av-kernel-group-law}
\emph{Step 1: lift $a$ along $A\to\Spec K$.}  By \Cref{prop:av-basic} the
structure morphism $A\to\Spec K$ is smooth, hence formally smooth, so there is
$\tilde u\colon T\to A$ over $K$ with $\tilde u|_{T_0}=a$.  (When the image of
$a$ lies in a single affine open $U\subseteq A$ --- which is the only case
needed for the affine-local reading of \Cref{rem:av-smooth-affine-local} ---
this is just the lifting property of the formally smooth $K$-algebra
$\Gamma(U,\O)$, and by \Cref{lem:av-lift-lands-in-open} the lift again lands in
$U$.)

\emph{Step 2: measure the error.}  Both $[n]\circ\tilde u$ and $b$ are elements
of $A(R)$, and they have the same restriction to $T_0$: indeed
$([n]\circ\tilde u)|_{T_0}=[n]\circ(\tilde u|_{T_0})=[n]\circ a=b|_{T_0}$.
Hence
\[
  c:=([n]\circ\tilde u)\cdot b^{-1}\in
  N(R,I)=\ker\bigl(A(R)\to A(R/I)\bigr).
\]

\emph{Step 3: correct the error.}  By \Cref{cor:av-kernel-npow} the map
$w\mapsto w^n$ is \emph{surjective} on $N(R,I)$, so there is $w\in N(R,I)$ with
$w^n=c^{-1}$.  Put $u:=\tilde u\cdot w$.  Then $u|_{T_0}=a\cdot1=a$, because
$w|_{T_0}=1$.  Moreover, using \Cref{lem:av-mulby-points} and the commutativity
of $A(R)$ (\Cref{prop:av-basic}),
\[
  [n]\circ u=u^n=\tilde u^{\,n}\cdot w^n=([n]\circ\tilde u)\cdot c^{-1}
  =([n]\circ\tilde u)\cdot\bigl(([n]\circ\tilde u)^{-1}\cdot b\bigr)=b .
\]
So $u$ is the required lift.
\end{proof}

\begin{remark}[Reading this affine-locally]\label{rem:av-smooth-affine-local}
\uses{thm:av-mulby-formally-smooth, lem:av-lift-lands-in-open}
The proof above is written with target $A$, but for the formalisation one wants
it for a pair of affine opens.  Let $U,V\subseteq A$ be affine opens with
$[n](U)\subseteq V$, let $B$ be a $\Gamma(V,\O)$-algebra, let $I\subseteq B$ be
a square-zero ideal and let $\Gamma(U,\O)\to B/I$ be a
$\Gamma(V,\O)$-algebra map.  Geometrically this is a morphism
$a\colon\Spec B/I\to U$ together with $b\colon\Spec B\to V$ as in the theorem.
Two observations make the affine-local statement follow from the proof above.
\begin{enumerate}
  \item In Step 1 the morphism $a$ lands in the affine open $U$, so only the
    lifting property of the $K$-algebra $\Gamma(U,\O)$ is needed --- that is,
    only formal smoothness of $\Gamma(U,\O)$ over $K$, which Mathlib provides
    from smoothness of $A\to\Spec K$ through \texttt{HasRingHomProperty}.  No
    formal smoothness for morphisms of schemes is required.
  \item The lift $u$ produced in Step 3 need not visibly land in $U$, but it
    does: $u|_{T_0}=a$ lands in $U$, so \Cref{lem:av-lift-lands-in-open}
    applies.
\end{enumerate}
Thus the whole of \Cref{sec:av-etale-unramified,sec:av-etale-smooth} can be
carried out with only ring-level formal smoothness and formal unramifiedness,
which is what Mathlib has.
\end{remark}

\subsection{$[n]$ is \'etale}\label{sec:av-etale-conclusion}

\begin{lemma}[$[n]$ is locally of finite presentation]\label{lem:av-mulby-lfp}
\uses{prop:av-basic, def:av-mulby}
\lean{AlgebraicGeometry.GrpObj.smooth_npow_left}
\leanok
The morphism $[n]\colon A\to A$ is locally of finite presentation.
\end{lemma}
\begin{proof}
\uses{prop:av-basic}
\leanok
By \Cref{prop:av-basic} the structure morphism $\pi\colon A\to\Spec K$ is
smooth, hence locally of finite presentation.  Now $[n]$ is a morphism of
$K$-schemes, so $\pi\circ[n]=\pi$; and a morphism $f\colon X\to Y$ of schemes
over $S$ with $X\to S$ and $Y\to S$ locally of finite presentation is itself
locally of finite presentation, by \stacksproject{02FV}.  Affine-locally this
is the statement that if $B$ and $C$ are finitely presented $K$-algebras then
$C$ is a finitely presented $B$-algebra for any $K$-algebra map $B\to C$, and
that is the form the formalisation uses: Mathlib \texttt{v4.33.1} has no
cancellation for \texttt{LocallyOfFinitePresentation} at the level of schemes,
so the Lean proof cancels at ring level with
\texttt{RingHom.FinitePresentation.of\_comp\_finiteType} applied to the
structure map $K\to\Gamma(V,\O)$ of an affine chart.
\end{proof}

\begin{theorem}[$[n]$ is \'etale for $n$ invertible]\label{thm:av-mulby-etale}
\uses{thm:av-mulby-formally-unramified, thm:av-mulby-formally-smooth,
lem:av-mulby-lfp}
\lean{AlgebraicGeometry.IsAbelianVariety.etale_mulBy,
AlgebraicGeometry.GrpObj.etale_npow_left}
\leanok
Let $A$ be an abelian variety over a field $K$ and let $n$ be invertible in
$K$.  Then $[n]\colon A\to A$ is \'etale.  This holds over an arbitrary field.
\end{theorem}
\begin{proof}
\uses{thm:av-mulby-formally-unramified, thm:av-mulby-formally-smooth,
lem:av-mulby-lfp}
\leanok
By \Cref{thm:av-mulby-formally-smooth} the morphism $[n]$ is formally smooth,
and by \Cref{lem:av-mulby-lfp} it is locally of finite presentation; a formally
smooth morphism that is locally of finite presentation is smooth
(\stacksproject{02H6}), and a smooth morphism is flat.  \emph{This is where
flatness comes from in the route written here: from
\Cref{thm:av-mulby-formally-smooth}, not from a dimension count and not from
miracle flatness.}  By \Cref{thm:av-mulby-formally-unramified} the morphism
$[n]$ is formally unramified.  Now
\[
  \text{Flat}\ +\ \text{FormallyUnramified}\ +\
  \text{LocallyOfFinitePresentation}\ \Longrightarrow\ \text{\'etale},
\]
which in Mathlib is
\texttt{AlgebraicGeometry.Etale.of\_formallyUnramified\_of\_flat}.
Equivalently one may use
\texttt{AlgebraicGeometry.Etale.eq\_smoothOfRelativeDimension\_zero}, which
identifies the \'etale morphisms with the smooth ones of relative dimension
$0$; on that route one has to produce the relative dimension $0$ instead, which
is no easier.
\end{proof}

\begin{remark}[The cheapest formalisation route]
\label{rem:av-etale-formalisation-route}
\uses{thm:av-mulby-etale, rem:av-smooth-affine-local,
rem:av-unramified-affine-local}
There is a shorter path that introduces no new class for schemes at all.  For
each pair of affine opens $U,V\subseteq A$ with $[n](U)\subseteq V$, prove
formal unramifiedness of $\Gamma(U,\O)$ over $\Gamma(V,\O)$ by
\Cref{thm:av-mulby-formally-unramified} read through
\Cref{rem:av-unramified-affine-local}, and formal smoothness of $\Gamma(U,\O)$
over $\Gamma(V,\O)$ by \Cref{thm:av-mulby-formally-smooth} read through
\Cref{rem:av-smooth-affine-local}.  Together these are
\texttt{Algebra.FormallyEtale}; with finite presentation
(\Cref{lem:av-mulby-lfp}) they are \texttt{Algebra.Etale}; and
\texttt{HasRingHomProperty} for \texttt{AlgebraicGeometry.Etale} and
\texttt{RingHom.Etale} promotes that to \'etaleness of $[n]$.  The
formalisation, in
\texttt{WeilConjectures/AlgebraicGeometry/Group/Etale.lean}, follows this
affine-local strategy but assembles it differently: it establishes the
scheme-level classes \texttt{AlgebraicGeometry.FormallyUnramified} and
\texttt{AlgebraicGeometry.Smooth} for $[n]$ chart by chart through
\texttt{HasRingHomProperty.iff\_appLE}, so that only
\texttt{Algebra.FormallyUnramified} and \texttt{Algebra.FormallySmooth} appear
at ring level, and concludes with
\texttt{AlgebraicGeometry.Etale.of\_formallyUnramified\_of\_flat} rather than
building \texttt{Algebra.Etale}.  Either way the mathematical content is
entirely \Cref{cor:av-kernel-npow}.
\end{remark}

\subsection{Corollaries for $n$ invertible}\label{sec:av-etale-corollaries}

\begin{theorem}[$[n]$ is finite, flat and surjective for $n$ invertible]
\label{thm:av-mulby-etale-corollaries}
\uses{thm:av-mulby-etale, thm:av-zariski-main, prop:av-basic}
\lean{AlgebraicGeometry.IsAbelianVariety.isFinite_mulBy_of_natCast_ne_zero,
AlgebraicGeometry.IsAbelianVariety.flat_mulBy_of_natCast_ne_zero,
AlgebraicGeometry.IsAbelianVariety.surjective_mulBy_of_natCast_ne_zero}
\leanok
Let $A$ be an abelian variety over a field $K$ and let $n$ be invertible in
$K$.  Then $[n]\colon A\to A$ is finite, flat and surjective.  This holds over
an arbitrary field.
\end{theorem}
\begin{proof}
\uses{thm:av-mulby-etale, thm:av-zariski-main, prop:av-basic, lem:av-mulby-lfp}
\leanok
\emph{Flat and locally quasi-finite.}  By \Cref{thm:av-mulby-etale} the
morphism $[n]$ is \'etale, hence flat, and hence locally quasi-finite: an
\'etale morphism is formally unramified and locally of finite type, and a
formally unramified morphism locally of finite type is locally quasi-finite.

\emph{Proper.}  Write $\pi\colon A\to\Spec K$ for the structure morphism, so
that $\pi\circ[n]=\pi$.  By \Cref{prop:av-basic}, $\pi$ is proper and
separated.  If $g\circ f$ is proper and $g$ is separated then $f$ is proper
(\stacksproject{01W6}); applying this with $f=[n]$ and $g=\pi$ shows that $[n]$
is proper.

\emph{Finite.}  By \Cref{thm:av-zariski-main} a proper, locally quasi-finite
morphism is finite; so $[n]$ is finite.

\emph{Surjective.}  The image $[n](A)$ is closed, because a finite morphism is
closed.  It is also open, because $[n]$ is flat and locally of finite
presentation (\Cref{lem:av-mulby-lfp}), and a flat morphism locally of finite
presentation is open (\stacksproject{01UA}).  Finally $[n](A)$ is non-empty and
$A$ is connected, both by \Cref{prop:av-basic}; a non-empty subset of a
connected space that is both open and closed is the whole space, so
$[n](A)=A$.
\end{proof}

\begin{corollary}[$A[n]$ is finite \'etale for $n$ invertible]
\label{cor:av-torsion-finite-etale}
\uses{thm:av-mulby-etale-corollaries, lem:av-torsion-pullback}
\lean{AlgebraicGeometry.IsAbelianVariety.isFinite_torsion_of_natCast_ne_zero,
AlgebraicGeometry.IsAbelianVariety.flat_torsion_of_natCast_ne_zero,
AlgebraicGeometry.IsAbelianVariety.etale_torsion}
\leanok
Let $A$ be an abelian variety over a field $K$ and let $n$ be invertible in
$K$.  Then $A[n]\to\Spec K$ is finite, flat and \'etale, of rank $\deg[n]$.
This holds over an arbitrary field.
\end{corollary}
\begin{proof}
\uses{thm:av-mulby-etale-corollaries, thm:av-mulby-etale,
lem:av-torsion-pullback, lem:av-degree-base-change}
\leanok
By \Cref{lem:av-torsion-pullback} the morphism $A[n]\to\Spec K$ is the base
change of $[n]$ along the identity section.  Finiteness, flatness and
\'etaleness are stable under base change and hold for $[n]$ by
\Cref{thm:av-mulby-etale-corollaries} and \Cref{thm:av-mulby-etale}; the rank
is preserved by \Cref{lem:av-degree-base-change}.
\end{proof}

\begin{remark}[Which statements need which field]\label{rem:av-field-hypotheses}
\uses{lem:av-translation, lem:av-translation-intertwines,
lem:av-translation-spreading, lem:av-kernel-cotangent, lem:av-kernel-group-law,
cor:av-kernel-npow, thm:av-mulby-formally-unramified,
thm:av-mulby-formally-smooth, thm:av-mulby-etale,
thm:av-mulby-etale-corollaries, cor:av-torsion-finite-etale,
prop:av-etale-torsion, prop:av-divisibility, prop:av-finite-etale-points}
This is where a blueprint earns its keep, because the Lean statements are all
over an arbitrary field while the classical accounts are all over an
algebraically closed one.
\begin{itemize}
  \item \textbf{Arbitrary field, no hypothesis at all.}
    \Cref{lem:av-translation} and \Cref{lem:av-translation-intertwines}, the
    existence of translations and the intertwining identity;
    \Cref{def:av-cotangent}, \Cref{lem:av-cotangent-affine},
    \Cref{lem:av-cotangent-product} and \Cref{lem:av-cotangent-group-law};
    \Cref{lem:av-kernel-cotangent}, \Cref{lem:av-kernel-group-law},
    \Cref{cor:av-kernel-module} and \Cref{cor:av-kernel-npow};
    \Cref{lem:av-lift-lands-in-open},
    \Cref{thm:av-mulby-formally-unramified},
    \Cref{thm:av-mulby-formally-smooth}, \Cref{lem:av-mulby-lfp},
    \Cref{thm:av-mulby-etale}, \Cref{thm:av-mulby-etale-corollaries} and
    \Cref{cor:av-torsion-finite-etale}.  In particular the \emph{whole} of the
    \'etale route is valid over an arbitrary field, and the ``$n$ prime to the
    characteristic'' of the classical statement of Proposition 2.8 is used only
    through ``$n$ is invertible in $K$'', which over a field is the same
    condition.
  \item \textbf{Algebraically closed.}  Only three statements need it.
    \Cref{lem:av-translation-spreading}, because it identifies the $K$-points
    with the closed points; \Cref{prop:av-divisibility}, for the same reason;
    and the second half of \Cref{prop:av-etale-torsion}, the count of
    $A[n](K)$ and the isomorphism $A[n](K)\cong(\Z/n)^{2g}$, because a finite
    \'etale $K$-scheme of rank $r$ has $r$ rational points only when $K$ is
    separably closed.  The first half of \Cref{prop:av-etale-torsion} --- that
    $A[n]$ is finite \'etale of rank $n^{2g}$ as a scheme --- needs no
    hypothesis on $K$.
  \item \textbf{Separably closed suffices.}  In
    \Cref{prop:av-finite-etale-points} the decomposition of a finite \'etale
    algebra as a product of copies of the base field needs only that $K$ is
    separably closed, and that is how Mathlib states it; the residue fields of
    a finite \'etale $K$-scheme are separable over $K$, so the passage from
    points of the scheme to $K$-points needs only separable closedness as well.
    We state the proposition over an algebraically closed field because that is
    the hypothesis carried by \Cref{prop:av-divisibility}.
  \item \textbf{Perfect.}  Nothing in this chapter needs $K$ perfect.  It is
    worth saying why, because Mathlib's tools for smoothness over a field ---
    \texttt{genericPoint\_mem\_smoothLocus\_of\_perfectField} and
    \texttt{dense\_smoothLocus\_of\_perfectField} --- are stated for morphisms
    to the spectrum of a perfect field, and would be the natural entry point
    for an argument through the smooth locus.  They do not apply to
    $[n]\colon A\to A$, whose target is not the spectrum of a field, and the
    route above avoids the smooth locus altogether.
\end{itemize}
\end{remark}

\section{The general case: the theorem of the cube}\label{sec:av-cube}

The \'etale route says nothing about the degree of $[n]$, and nothing at all
when $n$ is divisible by the characteristic of $K$.  Both come from the theory
of line bundles on abelian varieties.  This section states the five inputs that
are needed, with proof sketches and references, and derives the remaining parts
of \Cref{prop:av-isogeny} from them.  Two of the five, \Cref{thm:av-cube} and
\Cref{thm:av-mulby-pullback}, are proved in \Cref{chap:cube} and are quoted
here from there.  The other three are marked as open obligations: Mathlib
\texttt{v4.33.1} has no invertible sheaves on a scheme, no Picard group of a
scheme, no ampleness, no cohomology of coherent sheaves and no intersection
numbers, so \emph{none} of the five can currently even be stated in Lean, and
\Cref{sec:cube-cost} counts what the two proved ones would cost.  The
statements below carry full labels so that a later chapter can depend on
them.

\statusnote{Whether this project builds the layer of algebraic geometry that
the five inputs of this section require --- invertible sheaves and the Picard
group of a scheme, ampleness, cohomology of coherent sheaves, intersection
numbers --- is an open scope question, and it has not been decided.  The
alternative is to quote \Cref{prop:av-isogeny} as an external input, which is
the precedent set by \S2.5 of the curve estimates, where the representability
of $\mathrm{Pic}^0$ is marked ``Omitted --- quoted input''.  Until that question
is settled, the Lean declarations listed in \Cref{sec:av-status} as deliberate
placeholders stay unproved.}

\begin{theorem}[Theorem of the cube]\label{thm:av-cube}
Let $X$ and $Y$ be proper varieties over an algebraically closed field $K$, let
$Z$ be a connected reduced scheme of finite type over $K$, and let
$x_0\in X$, $y_0\in Y$, $z_0\in Z$ be rational points.  Let $L$ be a line
bundle on $X\times Y\times Z$ whose restrictions to
\[
  \{x_0\}\times Y\times Z,\qquad X\times\{y_0\}\times Z,\qquad
  X\times Y\times\{z_0\}
\]
are all trivial.  Then $L$ is trivial.
\end{theorem}
\begin{proof}
\uses{thm:cube-cube}
This is \Cref{thm:cube-cube}, word for word; it is proved in
\Cref{chap:cube} and is restated here so that \Cref{sec:av-cube} reads on its
own.  The proof there is the classical one of
\cite[\S6]{MumfordAbelianVarieties} and \cite[I, \S5]{MilneAbelianVarieties}:
the locus in $Z$ over which $L$ is trivial is closed by the seesaw theorem and
open by a K\"unneth argument on the square-zero thickenings of a point of it
together with the theorem on formal functions, hence is all of $Z$.
\end{proof}

\begin{theorem}[Pullback of a symmetric line bundle along $[n]$]
\label{thm:av-mulby-pullback}
\uses{def:av-mulby}
Let $A$ be an abelian variety over an algebraically closed field $K$ and let
$L$ be a line bundle on $A$.  Then for every $n\in\Z$
\[
  [n]^*L\;\cong\;L^{\otimes\frac{n^2+n}2}\otimes
  \bigl([-1]^*L\bigr)^{\otimes\frac{n^2-n}2}.
\]
In particular, if $L$ is \emph{symmetric}, that is $[-1]^*L\cong L$, then
$[n]^*L\cong L^{\otimes n^2}$.
\end{theorem}
\begin{proof}
\uses{thm:cube-mulby-pullback}
This is \Cref{thm:cube-mulby-pullback}, proved in \Cref{sec:cube-abelian};
what follows is a sketch of the argument given there, not a second proof, so
the statements it mentions are recorded as dependencies of
\Cref{thm:cube-mulby-pullback} in \Cref{chap:cube} and not again here.  Apply
\Cref{thm:av-cube} on
$A\times A\times A$ to the line bundle
\[
  m_{123}^*L\otimes m_{12}^*L^{-1}\otimes m_{13}^*L^{-1}\otimes
  m_{23}^*L^{-1}\otimes p_1^*L\otimes p_2^*L\otimes p_3^*L,
\]
where $m_S$ denotes the sum of the coordinates indexed by $S$ and $p_i$ the
projections: its restrictions to the three coordinate slices are trivial, so it
is trivial; that is \Cref{thm:cube-abelian}.  Evaluating the resulting identity
at the three points $[n]$, $[1]$ and $[-1]$ of the group $A(A)$ turns it into
the recursion
\[
  [n+1]^*L\otimes[n-1]^*L\cong\bigl([n]^*L\bigr)^{\otimes2}\otimes L\otimes
  [-1]^*L,
\]
of \Cref{cor:cube-mulby-recursion}, and the stated formula follows by a
two-sided induction on $n\in\Z$ from the three base cases $[0]^*L\cong\O_A$,
$[1]^*L\cong L$ and $[-1]^*L$.  See \cite[\S6]{MumfordAbelianVarieties} and
\cite[I, \S5]{MilneAbelianVarieties}.
\end{proof}

\begin{proposition}[Existence of a symmetric ample line bundle]
\label{prop:av-symmetric-ample}
Every abelian variety $A$ over an algebraically closed field carries an ample
line bundle $L$ with $[-1]^*L\cong L$.
\openobligation{AlgebraicGeometry: ampleness of an invertible sheaf on a
scheme, and the projectivity of an abelian variety.}
\end{proposition}
\begin{proof}
\emph{Sketch.}  An abelian variety is projective, so it carries some ample line
bundle $L_0$ \cite[\S6]{MilneAbelianVarieties}.  Then $L:=L_0\otimes[-1]^*L_0$
is ample, being a tensor product of ample bundles --- $[-1]$ is an isomorphism,
so $[-1]^*L_0$ is ample --- and it is symmetric because $[-1]$ is an
involution.
\end{proof}

\begin{proposition}[Positivity of the top self-intersection of an ample bundle]
\label{prop:av-ample-self-intersection}
Let $X$ be a proper variety of dimension $g$ over an algebraically closed field
and let $L$ be an ample line bundle on $X$.  Then the self-intersection number
$(L^g)$ is a positive integer.
\openobligation{AlgebraicGeometry: intersection numbers of invertible sheaves
on a proper scheme, via the Snapper polynomial and Euler characteristics of
coherent sheaves.}
\end{proposition}
\begin{proof}
\emph{Sketch.}  By Snapper's theorem the function
$m\mapsto\chi(X,L^{\otimes m})$ is a numerical polynomial of degree at most $g$
whose degree-$g$ coefficient is $(L^g)/g!$.  For $L$ ample and $m\gg0$ one has
$\chi(X,L^{\otimes m})=h^0(X,L^{\otimes m})$, which grows like a positive
multiple of $m^g$; so the leading coefficient, and with it $(L^g)$, is
positive.
\end{proof}

\begin{proposition}[Degree and intersection numbers]
\label{prop:av-degree-intersection}
\uses{def:av-degree}
Let $f\colon X\to Y$ be a finite, flat, surjective morphism of proper varieties
of the same dimension $g$ over an algebraically closed field, of degree $d$.
Then for every line bundle $L$ on $Y$
\[
  \bigl((f^*L)^g\bigr)=d\cdot\bigl(L^g\bigr).
\]
\openobligation{AlgebraicGeometry: intersection numbers of invertible sheaves
on a proper scheme, and the projection formula for a finite morphism.}
\end{proposition}
\begin{proof}
\emph{Sketch.}  By the projection formula
$f_*\bigl((f^*L)^{\otimes m}\bigr)\cong L^{\otimes m}\otimes f_*\O_X$, and
$f_*\O_X$ is locally free of rank $d$ because $f$ is finite and flat of degree
$d$.  Since $f$ is finite, cohomology is unchanged by $f_*$, so
$\chi\bigl(X,(f^*L)^{\otimes m}\bigr)=\chi\bigl(Y,L^{\otimes m}\otimes
f_*\O_X\bigr)$, whose Snapper polynomial has degree-$g$ coefficient
$d\cdot(L^g)/g!$.  Comparing leading coefficients gives the identity.  See
\cite[\S5]{MumfordAbelianVarieties}.
\end{proof}

\begin{theorem}[The degree of $[n]$]\label{thm:av-degree-formula}
\uses{thm:av-mulby-pullback, prop:av-symmetric-ample,
prop:av-ample-self-intersection, prop:av-degree-intersection, def:av-degree}
Let $A$ be an abelian variety of dimension $g$ over an algebraically closed
field $K$ and let $n\neq0$.  Suppose $[n]$ is finite, flat and surjective.
Then $\deg[n]=n^{2g}$.
\openobligation{AlgebraicGeometry: ampleness of an invertible sheaf on a
scheme and the projectivity of an abelian variety; intersection numbers of
invertible sheaves on a proper scheme, via the Snapper polynomial and Euler
characteristics of coherent sheaves; the projection formula for a finite
morphism; and everything the theorem of the cube needs, which is invertible
sheaves and the Picard group of a scheme together with the cohomology of
sheaves on a scheme in degrees zero and one and the direct images of a
coherent sheaf along a proper morphism, counted in the formalisation cost
section of the chapter on the theorem of the cube.}
\end{theorem}
\begin{proof}
\uses{thm:av-mulby-pullback, prop:av-symmetric-ample,
prop:av-ample-self-intersection, prop:av-degree-intersection}
Choose a symmetric ample $L$ by \Cref{prop:av-symmetric-ample}.  By
\Cref{thm:av-mulby-pullback}, $[n]^*L\cong L^{\otimes n^2}$, so
$\bigl(([n]^*L)^g\bigr)=n^{2g}(L^g)$ by multilinearity of the intersection
number.  On the other hand $[n]$ is finite, flat and surjective between proper
varieties of dimension $g$, so \Cref{prop:av-degree-intersection} gives
$\bigl(([n]^*L)^g\bigr)=\deg[n]\cdot(L^g)$.  By
\Cref{prop:av-ample-self-intersection} we have $(L^g)>0$, so it may be
cancelled, and $\deg[n]=n^{2g}$.
\end{proof}

\begin{theorem}[$[n]$ is an isogeny for every $n\neq0$]
\label{thm:av-isogeny-general}
\uses{thm:av-mulby-pullback, prop:av-symmetric-ample, thm:av-zariski-main,
prop:av-basic, lem:av-translation-intertwines}
Let $A$ be an abelian variety of dimension $g$ over an algebraically closed
field $K$ and let $n\neq0$ be a natural number, possibly divisible by the
characteristic of $K$.  Then $[n]$ is finite, flat and surjective.
\openobligation{AlgebraicGeometry: ampleness of an invertible sheaf on a
scheme and the projectivity of an abelian variety; everything the theorem of
the cube needs, which is invertible sheaves and the Picard group of a scheme
together with the cohomology of sheaves on a scheme in degrees zero and one
and the direct images of a coherent sheaf along a proper morphism, counted in
the formalisation cost section of the chapter on the theorem of the cube; and
in addition miracle flatness, that is, a finite surjective morphism between
regular schemes of the same dimension is flat, which Mathlib v4.33.1 does not
have in any form.}
\end{theorem}
\begin{proof}
\uses{thm:av-mulby-pullback, prop:av-symmetric-ample, thm:av-zariski-main,
prop:av-basic, lem:av-translation-intertwines, lem:av-torsion-pullback}
\emph{Sketch.}  Choose a symmetric ample $L$ by
\Cref{prop:av-symmetric-ample}.  By \Cref{thm:av-mulby-pullback},
$[n]^*L\cong L^{\otimes n^2}$, which is ample because $n\neq0$.  The
restriction of $[n]^*L$ to $A[n]$ is trivial, because $[n]$ maps $A[n]$ to the
identity point (\Cref{lem:av-torsion-pullback}).  So $A[n]$ is a proper scheme
carrying a line bundle that is at once ample and trivial, which forces
$\dim A[n]=0$; being proper over $K$ as well, $A[n]$ is finite.  By
\Cref{lem:av-translation-intertwines} every fibre of $[n]$ over a rational
point is a translate of $A[n]$, hence finite, so $[n]$ is locally
quasi-finite; it is proper by the cancellation argument in the proof of
\Cref{thm:av-mulby-etale-corollaries}; hence it is finite by
\Cref{thm:av-zariski-main}.  Surjectivity follows because the image is closed
of dimension $g=\dim A-\dim A[n]$ while $A$ is irreducible
(\Cref{prop:av-basic}).  Flatness follows from miracle flatness: $[n]$ is a
finite surjective morphism between smooth, hence regular, varieties of the same
dimension.  See \cite[\S6]{MumfordAbelianVarieties} and
\cite[\S7]{MilneAbelianVarieties}.
\end{proof}

\begin{remark}[Descending to an arbitrary field]
\label{rem:av-descent-to-any-field}
\uses{thm:av-isogeny-general, thm:av-degree-formula, def:av-abelian-variety}
\Cref{thm:av-degree-formula} and \Cref{thm:av-isogeny-general} are stated over
an algebraically closed field, as in \cite[\S6]{MumfordAbelianVarieties}, while
\Cref{prop:av-isogeny} is over an arbitrary field.  The passage is by descent:
an abelian variety over $K$ stays an abelian variety after base change along
$\Spec\bar K\to\Spec K$, the formation of $[n]$ commutes with that base change,
and finiteness, flatness, surjectivity and the degree of a morphism may all be
checked after the faithfully flat base change $\Spec\bar K\to\Spec K$.  The
stability of \Cref{def:av-abelian-variety} under base change along an extension
of the ground field is already an instance in
\texttt{WeilConjectures/Mathlib/AlgebraicGeometry/Group/Abelian.lean}.
\end{remark}

\section{Formalisation status}\label{sec:av-status}

The Lean sources are
\texttt{WeilConjectures/AlgebraicGeometry/AbelianVariety/Isogeny.lean},
\texttt{WeilConjectures/AlgebraicGeometry/AbelianVariety/Torsion.lean},
\texttt{WeilConjectures/AlgebraicGeometry/Group/Infinitesimal.lean} and
\texttt{WeilConjectures/AlgebraicGeometry/Group/Etale.lean}, together with the
mirror files under \texttt{WeilConjectures/Mathlib/}.

\emph{Formalised, with a complete Lean proof.}
\Cref{def:av-abelian-variety}; \Cref{prop:av-basic}; \Cref{def:av-mulby};
\Cref{lem:av-mulby-points}; \Cref{def:av-torsion} and
\Cref{lem:av-torsion-pullback}; \Cref{def:av-degree} and
\Cref{lem:av-degree-base-change}; \Cref{thm:av-zariski-main};
\Cref{prop:av-finite-etale-points}; \Cref{thm:av-finite-abelian-group};
\Cref{thm:av-mulby-formally-unramified},
\Cref{thm:av-mulby-formally-smooth} and \Cref{lem:av-mulby-lfp}, the last two
packaged into one declaration; \Cref{thm:av-mulby-etale};
\Cref{thm:av-mulby-etale-corollaries};
\Cref{cor:av-torsion-finite-etale}; and the two derivations
\Cref{prop:av-etale-torsion} and \Cref{prop:av-divisibility}.  The Lean proofs
of the last two are complete \emph{given} the four statements listed next,
which is exactly the sense in which a blueprint node can be green above a red
one.

\emph{Formalised as a deliberate \texttt{sorry}.}  Four declarations in
\texttt{WeilConjectures/AlgebraicGeometry/AbelianVariety/Isogeny.lean}:
\begin{itemize}
  \item \texttt{AlgebraicGeometry.IsAbelianVariety.isFinite\_mulBy} ---
    finiteness of $[n]$, part of \Cref{prop:av-isogeny};
  \item \texttt{AlgebraicGeometry.IsAbelianVariety.flat\_mulBy} --- flatness of
    $[n]$, part of \Cref{prop:av-isogeny};
  \item \texttt{AlgebraicGeometry.IsAbelianVariety.surjective\_mulBy} ---
    surjectivity of $[n]$, part of \Cref{prop:av-isogeny};
  \item \texttt{AlgebraicGeometry.IsAbelianVariety.finrank\_mulBy} ---
    $\deg[n]=n^{2g}$, that is \Cref{thm:av-degree-formula}.
\end{itemize}
All four are the content of \Cref{sec:av-cube}, are owned by issue 776, and
are blocked on the scope question recorded in the status note there.

The good-characteristic case, by contrast, is complete.  The first three of the
four are open only when $n$ is divisible by the characteristic of $K$:
\Cref{thm:av-mulby-etale-corollaries} gives finiteness, flatness and
surjectivity of $[n]$ for $n$ invertible in $K$ independently of
\Cref{sec:av-cube}, and those statements are formalised without a
\texttt{sorry} as \texttt{isFinite\_mulBy\_of\_natCast\_ne\_zero},
\texttt{flat\_mulBy\_of\_natCast\_ne\_zero} and
\texttt{surjective\_mulBy\_of\_natCast\_ne\_zero} in the same file, as is the
properness of $[n]$ that the first and third of them rest on (the general
\texttt{AlgebraicGeometry.isProper\_left} in
\texttt{WeilConjectures/Mathlib/AlgebraicGeometry/Morphisms/Proper.lean}).
Consequently \Cref{prop:av-divisibility} is proved with no \texttt{sorry} for
$n$ invertible in $K$ --- that is
\texttt{exists\_npow\_eq\_of\_natCast\_ne\_zero}, which like the node itself
still asks for $K$ algebraically closed --- and so is the finite-flat-\'etale
half of \Cref{prop:av-etale-torsion}, which is
\Cref{cor:av-torsion-finite-etale}.  What does \emph{not} follow for $n$
invertible is the rank $n^{2g}$, because \texttt{finrank\_mulBy} is open in
every characteristic; see \Cref{rem:av-etale-route-scope}.

\emph{Statement formalised, but proved by a different argument.}
\Cref{cor:av-kernel-npow}, as
\texttt{AlgebraicGeometry.GrpObj.eq\_one\_of\_pow\_eq\_one} and
\texttt{AlgebraicGeometry.GrpObj.exists\_pow\_eq}, both without a
\texttt{sorry}.  The Lean proof does not go through the cotangent space, so the
proof written here is not the formalised one and is left unmarked; the status
note at the head of \Cref{sec:av-etale-route} records the difference, and
writing the formalised argument out as a node of its own is a follow-up.

\emph{No Lean statement yet.}  Everything else in \Cref{sec:av-etale-route}:
the translations
\Cref{def:av-translation}, \Cref{lem:av-translation},
\Cref{lem:av-translation-intertwines} and
\Cref{lem:av-translation-spreading}; the cotangent-space material
\Cref{def:av-cotangent}, \Cref{lem:av-cotangent-affine},
\Cref{lem:av-cotangent-product} and \Cref{lem:av-cotangent-group-law}; the
kernel computation \Cref{lem:av-kernel-cotangent},
\Cref{lem:av-kernel-group-law} and \Cref{cor:av-kernel-module}; and
\Cref{lem:av-lift-lands-in-open}, which has no declaration of its own but whose
content the Lean proof does use, inline in the divisibility step of
\texttt{WeilConjectures/AlgebraicGeometry/Group/Etale.lean}.  And
everything in \Cref{sec:av-cube}:
\Cref{thm:av-cube}, \Cref{thm:av-mulby-pullback},
\Cref{prop:av-symmetric-ample}, \Cref{prop:av-ample-self-intersection},
\Cref{prop:av-degree-intersection}, \Cref{thm:av-degree-formula} and
\Cref{thm:av-isogeny-general} --- none of which can be stated in Mathlib
\texttt{v4.33.1} at all.
